\documentclass[10pt,a4paper]{amsart}
\usepackage[margin=19mm]{geometry}
\usepackage{amsmath,amssymb,mathtools}
\usepackage[scr=rsfs,cal=euler]{mathalfa}
\usepackage{xfrac}
\usepackage[unicode,hidelinks,bookmarks=false]{hyperref}
\newcommand{\ie}{i.e.\ }
\newcommand{\cf}{cf.\ }
\newcommand{\resp}{respectively\ }
\newcommand{\Subsection}{\subsection}
\providecommand{\nobreakdash}{\nobreak\hbox{-}}
\setlength{\emergencystretch}{2em}
\multlinegap0pt
\usepackage{bm}
\usepackage{bbm}
\usepackage{dsfont}
\usepackage{xspace}
\usepackage{cancel}
\usepackage{mathtools}
\usepackage{amsthm}
\makeatletter
\@ifundefined{theorem}{\newtheorem{theorem}{Theorem}}{}
\@ifundefined{proposition}{\newtheorem{proposition}[theorem]{Proposition}}{}
\makeatother
\newcommand{\qd}{\operatorname{qd}}
\newcommand{\score}{\operatorname{sc}}
\title[The maximum quartet distance between phylogenetic trees]{The maximum quartet distance between phylogenetic trees}
\author{Lior Pachter}
\date{Source: arXiv:2608.03542v1 (CC BY 4.0)}
\begin{document}
\maketitle

\noindent\emph{Source and licence.} The maximum quartet distance between phylogenetic trees. Version arXiv:2608.03542v1 (CC BY 4.0), distributed under the Creative Commons Attribution 4.0 International licence, as recorded in the arXiv licence metadata for this exact version and shown on the abstract page https://arxiv.org/abs/2608.03542v1. Full rights notice: licenses/arxiv-2608.03542-CC-BY-4.0.txt.\par\smallskip

\noindent\textit{Editorial scope.} Counted unit: the Proposition whose source label is \texttt{prop:elementary-caterpillar} in the article (source0.tex lines 691--708, source label \texttt{prop:elementary-caterpillar}), delivered together with the article's own complete original proof of it (source0.tex lines 710--766, one \texttt{proof} environment of 57 lines). No mathematics is written, repaired or re-derived: statement and proof are the article's own text line for line. Required context retained uncounted: prop:common-root-reduction (lines 515--526); eq:four-point-one-third (lines 687--689). Every definition, display and label of those lines is retained unchanged. Omitted from this package and not redistributed: 1--514, 527--686, 690--690, 709--709, 767--804. Nothing inside a retained excerpt is omitted or altered: every retained body is delivered exactly as the article's lines are written, so no omission receipt of any kind accompanies this package. Citations: the retained excerpts carry no citation command at all, so no bibliography entry of the article is transcribed and no reference is redistributed. Reference adaptation: none was needed and none was made; every reference inside the delivered unit resolves inside it, so no unresolved reference or citation is produced. Printed slips kept literal: no printed word, symbol, hypothesis or number of the counted statement or of its proof was altered. Layout and class: the article's own class is \texttt{article} with packages including fontenc, amsmath, amssymb, amsthm, graphicx, tikz, geometry, natbib, hyperref; the wrapper is the lane's pinned \texttt{amsart} editorial wrapper at 10pt on A4, which restates the theorem-like environments the retained text needs and adds the article's own macro definitions that the retained text needs (\texttt{\textbackslash qd}, \texttt{\textbackslash score}), with extra wrapper packages (bm, bbm, dsfont, xspace, cancel, mathtools). The delivered numbering is the unit's own. Assets: no retained span contains a figure environment, a graphics-inclusion command, a PDF-page inclusion, a table or a TikZ picture; no image file is copied into this package, the package declares no asset dependency and no image file is redistributed with it. Licence: version arXiv:2608.03542v1 of this article is distributed under the Creative Commons Attribution 4.0 International licence, as recorded in the arXiv licence metadata for this exact version and shown on the abstract page https://arxiv.org/abs/2608.03542v1. A full rights notice is delivered at licenses/arxiv-2608.03542-CC-BY-4.0.txt. Characters: every retained excerpt is pure ASCII, so the delivered engine, XeLaTeX with its default Latin Modern text font, reports no missing character.
\begin{proposition}[quantitative common-root reduction]
\label{prop:common-root-reduction}
Let $n\,\geq\,8$ and $E\,\geq\,0$.  Suppose every pair of ordered caterpillars on the
same $n$-element leaf set has score at least $-E$.  Then every pair
$T_1,T_2$ of binary phylogenetic trees on that leaf set satisfies
\[
 \qd(T_1,T_2)
 \,\leq\,\frac23\binom n4+
 \frac{4n}{3(n-4)}\left(E+2\binom{n-1}{3}\right)
 \,\leq\,\frac23\binom n4+\frac83(E+2n^3).
\]
\end{proposition}

\begin{equation}\label{eq:four-point-one-third}
 \mathbb P(\mathcal G)\,\geq\,\frac13.
\end{equation}

\begin{proposition}[elementary caterpillar bound]
\label{prop:elementary-caterpillar}
Every pair $C_1,C_2$ of ordered caterpillars on the same $n$-element leaf set,
with $n\,\geq\,4$, satisfies
\[
 \qd(C_1,C_2)
 \,\leq\,\left(\frac23+\frac6n\right)\binom n4
\]
and
\[
 \score(C_1,C_2)
 \,\geq\,-\frac{18}{n}\binom n4
 \,=\,-\frac34(n-1)(n-2)(n-3)
 \,\geq\,-\frac34n^3.
\]
The score bound is equivalent to the distance bound and is used in
Proposition~\ref{prop:common-root-reduction}.
\end{proposition}

\begin{proof}
Label the leaves in the order of $C_1$ as $1,\ldots,n$, and let $\pi(i)$ be
the position of leaf $i$ in the order of $C_2$; thus
$\pi\in S_n$.  For $i\,<\,j\,<\,k\,<\,\ell$, the two caterpillars agree
on this quartet precisely when
\[
 \max(\pi(i),\pi(j))\,<\,\min(\pi(k),\pi(\ell))
 \quad\text{or}\quad
 \max(\pi(k),\pi(\ell))\,<\,\min(\pi(i),\pi(j)).
\]
Call such a four-element subset good.

For a discrete application of \eqref{eq:four-point-one-third}, choose
$L$ uniformly from $\{1,\ldots,n\}$, let $U_1,U_2$ be independent uniform
random variables on $[0,1]$, independent also of $L$, and set
\[
 Z\,=\,\frac{L-1+U_1}{n},
 \qquad
 V\,=\,\frac{\pi(L)-1+U_2}{n}.
\]
Both marginals are uniform on $[0,1]$, the second because $\pi$ is a
permutation.  Take four independent copies and let $\mathcal L$ be the event
that their four values of $L$ are distinct.  Conditional on $\mathcal L$, the
set of those four values is uniform over
$\binom{\{1,\ldots,n\}}4$; their order by $Z$ is their order in $C_1$, while
their order by $V$ is their order in $C_2$.  Hence
\[
 \mathbb P(\mathcal G\mid\mathcal L)
 \,=\,\frac{A(\pi)}{\binom n4},
\]
where $A(\pi)$ is the number of good four-element subsets.  A union bound gives
\[
 \mathbb P(\mathcal L^{\mathsf c})
 \,\leq\,\binom42\frac1n\,=\,\frac6n.
\]
Together with \eqref{eq:four-point-one-third}, this implies
\[
 \frac13
 \,\leq\,\mathbb P(\mathcal G)
 \,\leq\,\mathbb P(\mathcal G\mid\mathcal L)
       +\mathbb P(\mathcal L^{\mathsf c})
 \,\leq\,\frac{A(\pi)}{\binom n4}+\frac6n.
\]
Therefore
\[
 A(\pi)\,\geq\,\left(\frac13-\frac6n\right)\binom n4,
\]
which gives the asserted distance bound because
$\qd(C_1,C_2)\,=\,\binom n4-A(\pi)$.  Finally, with
$N\,=\,\binom n4$,
\[
 \score(C_1,C_2)
 \,=\,3A(\pi)-N
 \,\geq\,-\frac{18}{n}\binom n4
 \,\geq\,-\frac34n^3.
\]
\end{proof}

\end{document}
